% =========================================================================
% Width-Bounded Heuristics for Symbolic Search
% Draft for ICAPS 2027. Compile with pdflatex; port to icaps27.sty later.
% Experiments section reports results from the SymK implementation.
% =========================================================================
\documentclass[letterpaper]{article}
\usepackage[margin=1in]{geometry}
\usepackage{amsmath,amssymb,amsthm}
\usepackage{natbib}
\usepackage{booktabs}
\usepackage{microtype}

% ------------------------------- macros ---------------------------------
\newcommand{\inlinecite}[1]{\citet{#1}}
\newcommand{\task}{\ensuremath{\Pi}}
\newcommand{\vars}{\ensuremath{\mathcal{V}}}
\newcommand{\ops}{\ensuremath{\mathcal{O}}}
\newcommand{\states}{\ensuremath{\mathcal{S}}}
\newcommand{\sinit}{\ensuremath{s_{\mathrm{I}}}}
\newcommand{\sgoal}{\ensuremath{s_\star}}
\newcommand{\pre}{\ensuremath{\mathit{pre}}}
\newcommand{\eff}{\ensuremath{\mathit{eff}}}
\newcommand{\cost}{\ensuremath{\mathit{cost}}}
\newcommand{\hstar}{\ensuremath{h^{*}}}
\newcommand{\hblind}{\ensuremath{h^{0}}}
\newcommand{\hpot}{\ensuremath{h^{\text{pot}}}}
\newcommand{\gstar}{\ensuremath{g^{*}}}
\newcommand{\copt}{\ensuremath{C^{*}}}
\newcommand{\bdd}[1]{\ensuremath{B_{#1}}}
\newcommand{\bddsize}[1]{\ensuremath{|B_{#1}|}}
\newcommand{\cofcount}{\ensuremath{Q}}
\newcommand{\width}{\ensuremath{\mathsf{w}}}
\newcommand{\values}{\ensuremath{V}}
\newcommand{\level}[2]{\ensuremath{H_{#2}}}
\newcommand{\effort}{\ensuremath{\mathit{effort}}}
\newcommand{\sub}[2]{\ensuremath{\mathrm{cof}_{#1}(#2)}}
\newcommand{\img}{\ensuremath{\mathrm{img}}}
\newcommand{\bddastar}{BDDA\ensuremath{^{*}}}
\newcommand{\astar}{A\ensuremath{^{*}}}
\newcommand{\ghseta}{GHSETA\ensuremath{^{*}}}
\newcommand{\symba}{SymBA\ensuremath{^{*}}}

\theoremstyle{plain}
\newtheorem{theorem}{Theorem}
\newtheorem{lemma}{Lemma}
\newtheorem{proposition}{Proposition}
\newtheorem{corollary}{Corollary}
\theoremstyle{definition}
\newtheorem{definition}{Definition}
\newtheorem{example}{Example}
\theoremstyle{remark}
\newtheorem{remark}{Remark}

\title{Heuristics with Bounded Cofactor Width for Symbolic Search}
\author{Anonymous for review}
\date{}

\begin{document}
\maketitle

\begin{abstract}
Symbolic search with binary decision diagrams (BDDs) is a leading
approach to cost-optimal classical planning, but its dominant
instantiation remains blind: partitioning the search frontier by
heuristic value can increase BDD sizes so much that even the perfect
heuristic deteriorates performance exponentially. We give the matching
sufficient condition for safety. A heuristic has \emph{cofactor width}
$W$ under the search's variable order if its value is computable by
scanning the state variables once in that order while distinguishing at
most $W$ cases. We prove that partitioning \emph{any} state set by such
a heuristic inflates the total BDD size by at most a factor polynomial
in $W$ and the number of variables, that symbolic \astar\ with a
consistent width-$W$ heuristic incurs at most this polynomial overhead
over blind forward search, and that the known exponential deterioration
requires exponential width. Integer potential heuristics, prefix pattern
databases and linear merge-and-shrink abstractions instantiate the class
with widths controlled by their standard parameters, yielding the first
formal performance guarantee for the operator-potential approach.
Experiments in the SymK planner confirm the theory: fragmentation stays
negligible across five orders of magnitude of measured width, the
potential-magnitude cap acts as a knob with bounded overhead, and the
merge-and-shrink family surpasses blind forward search---while an
alignment study shows that width is a currency to be spent on accuracy,
not an objective in itself.
\end{abstract}

\section{Introduction}

A common approach to solving classical planning tasks optimally is
\astar\ search \citep{hart-et-al-ieeessc1968} with an admissible heuristic
\citep{pearl-1984}. Symbolic search
\citep{mcmillan-1993,edelkamp-helmert-ecp1999} replaces the expansion of
individual states with the expansion of state sets, compactly represented
as binary decision diagrams \citep[BDDs;][]{bryant-ieeecomp1986}. Symbolic
bidirectional blind search is among the strongest algorithms for
cost-optimal planning
\citep{torralba-et-al-aij2017,speck-et-al-jair2025}, and a key component
of the winning entry of the optimal track of the International Planning
Competition (IPC) 2023.

Combining the two paradigms has a long history. \bddastar\
\citep{edelkamp-reffel-ki1998} represents a heuristic $h$ as one BDD per
heuristic value and splits each search layer into buckets of states that
share a $g$- and an $h$-value. SetA\ensuremath{^{*}} and \ghseta\
\citep{jensen-et-al-aij2008} instead fold the heuristic-value change into
the transition relations, and ADDA\ensuremath{^{*}}
\citep{hansen-et-al-sara2002} stores $h$ as an algebraic decision diagram
(ADD). Empirically, all of these variants beat blind symbolic search in
some domains and lose in others
\citep{torralba-phd2015,torralba-et-al-aij2017}.
\citet{speck-et-al-icaps2020} explained why: in symbolic search the
effort of an expansion is governed by the number of BDD nodes rather than
the number of represented states, and they prove that even the perfect
heuristic \hstar\ can increase the cumulative size of the expanded BDDs
exponentially. Consequently, heuristic symbolic search lost traction, the
one notable exception being operator-potential heuristics
\citep{fiser-et-al-aaai2022,fiser-et-al-aij2024}, which integrate
LP-based potential heuristics \citep{pommerening-et-al-aaai2015} into
\ghseta\ and empirically ``hardly ever suffer'' from increased BDD sizes.
Why this particular family behaves so well has remained without formal
explanation, and no condition is known under which a heuristic
\emph{provably} cannot hurt symbolic search. Recently,
\citet{speck-helmert-ecai2025} gave the first worst-case performance
guarantees for symbolic search relative to explicit search, but their
analysis concerns the representation of transition relations and goal
formulas and explicitly leaves heuristics aside.

In this paper we close this gap. We introduce \emph{width-bounded
heuristics}: a heuristic has cofactor width $W$ under a variable order
if its value can be determined by reading the state variables once in
that order while maintaining at most $W$ distinct cases at any point,
i.e., if its algebraic decision diagram in the search's variable order
has at most $W$ nodes per level. Our contributions are as follows.
First, we prove a \emph{fragmentation theorem}: partitioning an arbitrary
state-set BDD by the values of a width-$W$ heuristic yields buckets whose
total size is bounded by a polynomial in $W$ times the size of the
original BDD---for every state set, every layer and every task. The bound
covers both known failure modes of heuristic symbolic search, the
fragmentation of the frontier and the fact that a subset of a state set
can have an exponentially larger BDD. Second, we prove a
\emph{search-effort theorem}: \bddastar\ with a consistent width-$W$
heuristic expands BDDs whose cumulative size exceeds that of blind
symbolic forward search by at most a factor $O(W^{2}n)$, where $n$ is
the number of binary variables. Third, we show that the pathological
heuristic in the construction of \citet{speck-et-al-icaps2020} has width
$2^{\Omega(n)}$, so exponential width is necessary for their exponential
deterioration. Fourth, we identify natural width-bounded families:
integer potential heuristics, pattern databases whose pattern is a prefix
of the variable order, and merge-and-shrink heuristics with a linear
merge strategy---whose shrink size limit \emph{is} a width bound. This
yields, to our knowledge, the first formal performance guarantee for the
operator-potential approach.
Finally, we implement all of the above in the SymK planner
\citep{speck-et-al-jair2025} and evaluate it on the IPC benchmarks: the
fragmentation predicted harmless stays harmless, the width knob trades
accuracy against a bounded and empirically small overhead, and the
merge-and-shrink family surpasses blind forward search.

% ============================ background =================================
\section{Background}

\subsection{Classical Planning}

We consider planning tasks in the SAS$^{+}$ formalism
\citep{backstrom-nebel-compint1995}.

\begin{definition}[planning task]
A \emph{planning task} is a tuple
$\task = \langle \vars, \ops, \sinit, \sgoal \rangle$, where \vars\ is a
finite set of \emph{state variables} $v$, each with a finite domain
$D_v$, \ops\ is a finite set of \emph{operators}
$o = \langle \pre_o, \eff_o \rangle$ given by partial variable
assignments and each associated with a cost
$\cost(o) \in \mathbb{Z}_{\geq 1}$, \sinit\ is a state (a total
assignment) and \sgoal\ is a partial assignment.
\end{definition}

Operator $o$ is applicable in state $s$ iff $s \models \pre_o$, and
applying it yields the state $s[o]$ that agrees with $\eff_o$ where
defined and with $s$ elsewhere. A \emph{plan} for $s$ is a sequence of
operators leading from $s$ to a goal state; its cost is the sum of the
operator costs. The perfect heuristic $\hstar(s)$ is the cost of a
cheapest plan for $s$, or $\infty$ if none exists, and
$\gstar(s)$ is the cost of a cheapest path from \sinit\ to $s$. We write
$\copt = \hstar(\sinit)$ and assume $0 < \copt < \infty$. A
\emph{heuristic} is a function
$h \colon \states \to \mathbb{Z}_{\geq 0} \cup \{\infty\}$. It is
\emph{admissible} if $h \leq \hstar$ and \emph{consistent} if
$h(s) \leq h(s[o]) + \cost(o)$ for all $s$ and applicable $o$, and
$h(s) = 0$ for all goal states. The blind heuristic is
$\hblind \equiv 0$. Restricting costs to positive integers excludes
zero-cost operators; this keeps the layer structure of uniform-cost
search simple and follows \citet{speck-et-al-icaps2020}, who assume unit
costs.

\subsection{Binary Decision Diagrams}

We fix a binary encoding of states: each variable $v$ is encoded by
$\lceil \log_2 |D_v| \rceil$ Boolean variables, and we assume every
Boolean assignment that does not encode a state is mapped to a designated
non-state value (Remark~\ref{rem-encoding} discusses this choice). Let
$x_1, \dots, x_n$ denote the Boolean variables in a fixed order $\pi$; we
call $\pi$ \emph{variable-contiguous} if the bits of each SAS$^{+}$
variable are consecutive. A set of states $S$ is represented by its
characteristic function $\chi_S \colon \{0,1\}^n \to \{0,1\}$ as a
reduced ordered BDD $\bdd{S}$ with respect to $\pi$
\citep{bryant-ieeecomp1986}. We write $\bddsize{S}$ for the number of inner
nodes of $\bdd{S}$. For a Boolean or integer-valued function $f$ over
$x_1, \dots, x_n$ and a prefix assignment $\rho \in \{0,1\}^i$ to
$x_1, \dots, x_i$, the \emph{cofactor} $f|_\rho$ is the function of
$x_{i+1}, \dots, x_n$ obtained by fixing the prefix. We write
$\sub{i}{f} = \{\, f|_\rho \mid \rho \in \{0,1\}^i \,\}$ for the set of
distinct cofactors of $f$ after $i$ variables, with
$\sub{0}{f} = \{f\}$. Reduced ordered BDDs are canonical: the nodes of
$\bdd{f}$ labeled $x_{i+1}$ correspond exactly to the cofactors in
$\sub{i}{f}$ that essentially depend on $x_{i+1}$
\citep{wegener-2000}. The following quantity, which counts all cofactors
including those represented implicitly by long edges and constants, is
central to our analysis.

\begin{definition}[cofactor count]
\label{def-cofcount}
For a function $f$ over $x_1, \dots, x_n$, the \emph{cofactor count} is
$\cofcount(f) = \sum_{i=0}^{n-1} |\sub{i}{f}|$.
\end{definition}

$\cofcount(f)$ equals the number of inner nodes of the quasi-reduced
(complete) OBDD of $f$, in which every root-to-sink path tests every
variable. Both size measures are polynomially related:

\begin{lemma}
\label{lem-qsize}
For every Boolean $f$ over $x_1, \dots, x_n$,
$\bddsize{f} \leq \cofcount(f) \leq n \cdot (\bddsize{f} + 2)$.
\end{lemma}

\begin{proof}
For the first inequality, consider a node $N$ of $\bdd{f}$ labeled
$x_{i+1}$. Since every node of a reduced OBDD is reachable from the root,
there is a path from the root to $N$; fixing the variables tested along
this path to the corresponding edge values and all remaining variables
among $x_1, \dots, x_i$ arbitrarily yields a prefix $\rho \in \{0,1\}^i$
with $f|_\rho$ equal to the function computed by $N$, because variables
skipped by long edges do not influence the traversal. Hence every node
labeled $x_{i+1}$ computes an element of $\sub{i}{f}$, and by canonicity
distinct nodes compute distinct functions, so the map is injective.
Summing over $i$ gives
$\bddsize{f} \leq \sum_{i=0}^{n-1} |\sub{i}{f}| = \cofcount(f)$. For the
second inequality, fix $i$ and consider any $u \in \sub{i}{f}$. Either
$u$ is a constant function, of which there are at most two, or $u$
essentially depends on some variable; let $x_j$ with $j > i$ be the first
such variable. Traversing $\bdd{f}$ from the root under any prefix
witnessing $u$ ends at an inner node labeled $x_j$ that computes $u$, and
by canonicity this node is unique. Hence
$|\sub{i}{f}| \leq \bddsize{f} + 2$ for every $i$, and summing over the
$n$ levels yields the claim.
\end{proof}

The factor $n$ in Lemma~\ref{lem-qsize} is unavoidable: for
$f = x_1$ we have $\bddsize{f} = 1$ but $\cofcount(f) = 2n - 1$, since
the constants $0$ and $1$ are cofactors at every level $i \geq 1$.

\subsection{Symbolic Heuristic Search and Expansion Size}

Symbolic uniform-cost search maintains, for increasing $g$, the BDD of
the layer $L_g = \{\, s \mid \gstar(s) = g \,\}$, obtained by an image
operation on the previous layers followed by removal of the closed
states \citep{torralba-et-al-aij2017}. \bddastar\
\citep{edelkamp-reffel-ki1998} additionally represents the heuristic $h$
by one BDD per value: for $v \in \img(h)$, the \emph{level set} is
$\level{h}{v} = \{\, s \mid h(s) = v \,\}$. Successor sets are
intersected with the level sets, and the resulting \emph{buckets} of
states sharing a $g$-value and an $h$-value are expanded in order of
increasing $f = g + v$. Following
\citet{kissmann-edelkamp-aaai2011,torralba-phd2015,speck-et-al-icaps2020},
ties are broken in favor of smaller $g$-values, all states with the same
$(g, v)$ are represented by a single BDD, generated states already in the
closed list are removed before bucketing, states with $h$-value $\infty$
are discarded, and the search terminates as soon as a selected bucket
contains a goal state.\footnote{Fixing the goal test to selected rather
than generated buckets matches the reading of
\citet{speck-et-al-icaps2020} that all buckets counted by
Definition~\ref{def-effort} \emph{must} be expanded before an optimal
plan is returned; they note that for their task family the choice is
immaterial.} We measure search effort exactly as
\citet{speck-et-al-icaps2020}:

\begin{definition}[expansion size, \citealp{speck-et-al-icaps2020}]
\label{def-effort}
Let $B_{\task,h}(g,v)$ denote the unique BDD expanded by \bddastar\ with
$g$-value $g$ and $h$-value $v$ on task \task\ with heuristic $h$, with
$B_{\task,h}(g,v)$ representing the empty set if no such expansion
occurs. The \emph{expansion size} of \bddastar\ is
$\effort(\task,h) = \sum_{\,g + v \leq \copt,\ g < \copt}
|B_{\task,h}(g,v)|$.
\end{definition}

For $\hblind$ all states have $h$-value $0$, and \bddastar\ coincides
with symbolic uniform-cost search, so
$\effort(\task,\hblind) = \sum_{g < \copt} \bddsize{L_g}$.
\citet{speck-et-al-icaps2020} prove that no analogue of the classical
domination result for consistent heuristics in explicit \astar\ holds:
there is a family of tasks $\task_n$ on which
$\effort(\task_n, \hstar)$ is exponentially larger than
$\effort(\task_n, \hblind)$. The reason is that a strict subset
$S' \subsetneq S$ can have an exponentially larger BDD than $S$
\citep{torralba-phd2015,speck-et-al-icaps2020}, so both splitting a
layer into buckets and pruning states from it can increase
representation size. Our goal is a property of $h$ that rules out both
effects for every state set at once.

% ===================== width-bounded heuristics ==========================
\section{Bounded Cofactor Width}

Our key definition measures how much information about the prefix of a
state a heuristic must retain in order to determine its value from the
remaining suffix. We treat a heuristic as a function
$h \colon \{0,1\}^n \to \mathbb{Z}_{\geq 0} \cup \{\infty\}$ over the
Boolean encoding.

\begin{definition}[cofactor width]
\label{def-width}
The \emph{cofactor width} of a heuristic $h$ under variable order $\pi$
is $\width_\pi(h) = \max_{0 \leq i \leq n} |\sub{i}{h}|$. We call $h$
\emph{$(\pi,W)$-width-bounded} if $\width_\pi(h) \leq W$, and write
\emph{width} for cofactor width throughout.
\end{definition}

Equivalently, $h$ is $(\pi,W)$-width-bounded iff it is computed by a
deterministic layered automaton that reads $x_1, \dots, x_n$ in the order
$\pi$ using at most $W$ states per level: the automaton states at level
$i$ are exactly the cofactors in $\sub{i}{h}$, the transition on reading
a value $b$ for $x_{i+1}$ maps the cofactor $u$ to $u|_{x_{i+1} = b}$,
and the output of a final-level state, which is a constant function, is
its value. In OBDD terminology, $\width_\pi(h)$ is the maximal number of
nodes per level of the quasi-reduced multi-terminal BDD of $h$ under
$\pi$---the classical width of a Boolean or pseudo-Boolean function
\citep{wegener-2000}, here applied to heuristics and relativized to the
order used by the search.

A word on terminology. \emph{Cofactor width} is a property of the
heuristic's \emph{representation} under a variable order and is
unrelated to the novelty width of \citet{lipovetzky-geffner-ecai2012},
which measures the size of value tuples needed to distinguish states
during search, and only indirectly related to structural graph
parameters such as pathwidth, whose boundedness is a classical
sufficient condition for small OBDDs \citep{wegener-2000}. Three
immediate consequences will do much of the work below.

\begin{proposition}[value range]
\label{prop-range}
If $\width_\pi(h) \leq W$, then $|\img(h)| \leq W$.
\end{proposition}

\begin{proof}
Every cofactor in $\sub{n}{h}$ is a constant function, one per attained
value, so $|\img(h)| = |\sub{n}{h}| \leq \width_\pi(h) \leq W$.
\end{proof}

We write $\values = |\img(h) \setminus \{\infty\}|$ for the number of
finite values of $h$; by Proposition~\ref{prop-range},
$\values \leq W$.

\begin{proposition}[closure]
\label{prop-closure}
Let $h_1, h_2$ have widths $W_1, W_2$ under $\pi$ and let
$\oplus \in \{+, \max, \min\}$. Then
$\width_\pi(h_1 \oplus h_2) \leq W_1 \cdot W_2$.
\end{proposition}

\begin{proof}
For every prefix $\rho$ we have
$(h_1 \oplus h_2)|_\rho = h_1|_\rho \oplus h_2|_\rho$, so the map
$(h_1|_\rho, h_2|_\rho) \mapsto (h_1 \oplus h_2)|_\rho$ is a surjection
from a set of size at most $W_1 W_2$ onto $\sub{i}{h_1 \oplus h_2}$ for
every $i$.
\end{proof}

Hence maximization and admissible additive combinations of few
width-bounded heuristics remain width-bounded, at a multiplicative cost
per component. Finally, width-boundedness is essentially equivalent to
the compactness of the heuristic's ADD in the search's variable order.
Let $\mathrm{ADD}_\pi(h)$ denote the reduced ordered ADD of $h$ under
$\pi$ \citep{bahar-et-al-iccad1993} and $|\mathrm{ADD}_\pi(h)|$ its
number of inner nodes.

\begin{proposition}[ADD characterization]
\label{prop-add}
$\width_\pi(h) \leq |\mathrm{ADD}_\pi(h)| + |\img(h)|$ and
$|\mathrm{ADD}_\pi(h)| \leq n \cdot \width_\pi(h)$.
\end{proposition}

\begin{proof}
For the first inequality, every cofactor in $\sub{i}{h}$ is either a
constant, of which there are $|\img(h)|$ many, or the function computed
by a unique inner node of $\mathrm{ADD}_\pi(h)$ at a level below $i$, by
canonicity of reduced ADDs. For the second, every inner node of
$\mathrm{ADD}_\pi(h)$ labeled $x_{i+1}$ computes a distinct cofactor in
$\sub{i}{h}$, so the nodes per level number at most $\width_\pi(h)$, and
there are $n$ levels.
\end{proof}

Width-bounded heuristics are thus exactly the heuristics with small ADDs
\emph{in the variable order used by the search}. This order dependence is
essential and is what separates our results from statements about the
existence of some good order: the search fixes one $\pi$ for state sets,
transition relations and heuristic alike
\citep{kissmann-hoffmann-jair2014}.

\begin{remark}
\label{rem-encoding}
Definitions over $\{0,1\}^n$ require fixing $h$ on Boolean assignments
that encode no state. All upper bounds below hold for whichever total
extension is used to build the level-set BDDs, so an implementation
should pick an extension of low width; for the families in
Section~\ref{sec-instances} the natural arithmetic or table-based
extension works. Lower bounds only ever evaluate $h$ on genuine states
and therefore hold for every extension.
\end{remark}

% ======================= fragmentation theorems ==========================
\section{Bounded Fragmentation}
\label{sec-fragmentation}

We now bound the representational cost of splitting an arbitrary state
set by heuristic values. Throughout, fix the order $\pi$ and write
$W = \width_\pi(h)$. Recall the classical product bound
$\bddsize{f \wedge g} \leq \bddsize{f} \cdot \bddsize{g}$
\citep{bryant-ieeecomp1986}. Applied naively with $g = \chi_{\level{h}{v}}$ it
bounds a single bucket by
$\bddsize{S} \cdot \bddsize{\level{h}{v}}$, which does not control the
sum over all buckets. The next two theorems give bounds whose sum over
buckets stays small. The proofs use one observation, which we isolate.

\begin{lemma}
\label{lem-pair}
Let $S \subseteq \{0,1\}^n$, let $U \subseteq \img(h)$ and let
$\rho, \rho' \in \{0,1\}^i$ be prefixes with
$\chi_S|_\rho = \chi_S|_{\rho'}$ and $h|_\rho = h|_{\rho'}$. Then
$\chi_{S \cap h^{-1}(U)}|_\rho = \chi_{S \cap h^{-1}(U)}|_{\rho'}$.
\end{lemma}

\begin{proof}
For every suffix $\sigma \in \{0,1\}^{n-i}$,
$\chi_{S \cap h^{-1}(U)}|_\rho(\sigma) = \chi_S|_\rho(\sigma) \cdot
[h|_\rho(\sigma) \in U]$, which by assumption equals
$\chi_S|_{\rho'}(\sigma) \cdot [h|_{\rho'}(\sigma) \in U]
= \chi_{S \cap h^{-1}(U)}|_{\rho'}(\sigma)$.
\end{proof}

In other words, the cofactor of a bucket is determined by the pair
consisting of the cofactor of the state set and the cofactor of the
heuristic. Counting realized pairs yields our first main result.

\begin{theorem}[bucket bound]
\label{thm-bucket}
Let $S \subseteq \{0,1\}^n$ and $U \subseteq \img(h)$. Then
\[
\bddsize{S \cap h^{-1}(U)}
\;\leq\; \cofcount(\chi_{S \cap h^{-1}(U)})
\;\leq\; W \cdot \cofcount(\chi_S)
\;\leq\; W n \, (\bddsize{S} + 2).
\]
\end{theorem}

\begin{proof}
The first and last inequalities are Lemma~\ref{lem-qsize}. For the
middle inequality, fix $i \in \{0, \dots, n-1\}$ and define the set of
\emph{realized pairs}
$R_i = \{\, (\chi_S|_\rho,\, h|_\rho) \mid \rho \in \{0,1\}^i \,\}$.
Consider the map
$\mu \colon \sub{i}{\chi_{S \cap h^{-1}(U)}} \to R_i$ defined as
follows: for $u \in \sub{i}{\chi_{S \cap h^{-1}(U)}}$, choose any
prefix $\rho$ with $u = \chi_{S \cap h^{-1}(U)}|_\rho$ and set
$\mu(u) = (\chi_S|_\rho, h|_\rho)$. The map is well defined on the
chosen representatives, and it is injective: if
$\mu(u) = \mu(u')$ via prefixes $\rho, \rho'$, then
$\chi_S|_\rho = \chi_S|_{\rho'}$ and $h|_\rho = h|_{\rho'}$, so
$u = u'$ by Lemma~\ref{lem-pair}. Hence
$|\sub{i}{\chi_{S \cap h^{-1}(U)}}| \leq |R_i|
\leq |\sub{i}{\chi_S}| \cdot |\sub{i}{h}|
\leq W \cdot |\sub{i}{\chi_S}|$.
Summing over $i = 0, \dots, n-1$ gives
$\cofcount(\chi_{S \cap h^{-1}(U)}) \leq W \cdot \cofcount(\chi_S)$.
\end{proof}

With $U = \{v\}$, Theorem~\ref{thm-bucket} bounds the size of every
single bucket $S \cap \level{h}{v}$. Beyond that, it bounds the size of
\emph{every} state set obtained by removing value slices from $S$: with
$U = \img(h) \setminus \{\infty\}$ it caps the non-dead-end subset of
$S$, and with $U = \{\, v \mid v \leq c \,\}$ it caps $f$-value-based
pruning. It therefore also controls the subset blowup that drives the
negative result of \citet{speck-et-al-icaps2020}: no subset of $S$
selected by heuristic values can outgrow $S$ by more than a factor
$Wn$, up to additive terms. Summing over all values and applying
Proposition~\ref{prop-range} bounds the entire partition.

\begin{theorem}[partition bound]
\label{thm-partition}
Let $S \subseteq \{0,1\}^n$. Then
\[
\sum_{v \in \img(h)} \bddsize{S \cap \level{h}{v}}
\;\leq\; W^{2} \cdot \cofcount(\chi_S)
\;\leq\; W^{2} n \, (\bddsize{S} + 2).
\]
\end{theorem}

\begin{proof}
By Theorem~\ref{thm-bucket} applied to $U = \{v\}$, each of the at most
$|\img(h)| \leq W$ (Proposition~\ref{prop-range}) buckets has BDD size
at most $W \cdot \cofcount(\chi_S)$; the second inequality is
Lemma~\ref{lem-qsize}.
\end{proof}

\begin{remark}
The factor $W$ contributed by Proposition~\ref{prop-range} can be
refined: in the proof of Theorem~\ref{thm-bucket}, the pair
$(\chi_S|_\rho, h|_\rho)$ contributes a nonconstant cofactor only for
values $v \in \img(h|_\rho)$, so the sum over $v$ of the cofactor counts
is bounded by
$\sum_{i} \sum_{(u,c) \in R_i} |\img(c)|$, where $c$ ranges over
heuristic cofactors. The number of values attainable from a cofactor is
often far smaller than $W$; implementations that share isomorphic
subgraphs across buckets in a unique table, as CUDD does, profit
accordingly.
\end{remark}

\begin{example}
\label{ex-parity}
The factor $n$ in Theorems~\ref{thm-bucket} and~\ref{thm-partition}
stems from Lemma~\ref{lem-qsize} and cannot be removed when sizes are
measured in reduced BDDs. Let $S = \{\, s \mid s(x_1) = 1 \,\}$, so
$\bddsize{S} = 1$, and let $h(s) = \bigl(\sum_{i=1}^{n} s(x_i)\bigr)
\bmod 2$, which has width $2$ under any order. The bucket
$S \cap \level{h}{0}$ consists of the assignments with $x_1 = 1$ and an
odd number of ones among $x_2, \dots, x_n$; its reduced BDD consists of
one node for $x_1$ and a parity chain with one node for $x_2$ and two
nodes for each of $x_3, \dots, x_n$, i.e., $2n - 2$ nodes in total: a
factor $2n - 2$ larger than $\bdd{S}$ although $W = 2$. Measured in
cofactor counts, where $\cofcount(\chi_S) = 2n - 1$, the bound of
Theorem~\ref{thm-bucket} is met within a factor of about two.
\end{example}

% ======================= search-effort guarantees ========================
\section{Search-Effort Guarantees}
\label{sec-effort}

We lift the set-level bounds to a guarantee on the expansion size of
\bddastar. The first step identifies the expanded buckets with slices of
the blind layers. This correspondence is folklore for explicit \astar\
with consistent heuristics; we prove it in the symbolic bucket setting
because the completeness of buckets at expansion time, i.e., that a
bucket is never re-populated after its expansion, is what makes
Definition~\ref{def-effort} well defined and our accounting exact.

\begin{lemma}[bucket identity]
\label{lem-identity}
Let $h$ be consistent. \bddastar\ on \task\ with $h$ selects every
bucket key $(g,v)$ with $g + v \leq \copt$ and $g < \copt$ and
$L_g \cap \level{h}{v} \neq \emptyset$ exactly once before terminating,
and the BDD expanded with key $(g,v)$ represents exactly
$L_g \cap \level{h}{v}$. No other bucket is expanded before
termination.
\end{lemma}

\begin{proof}
Throughout, the \emph{key} of a generated state $s$ is the pair
$(c, h(s))$, where $c$ is the cost of the generating path, and the
$f$-value of a key $(g,v)$ is $g + v$. We prove five claims.

\emph{Claim 1 (monotone selection): if a key $(g_1,v_1)$ is selected at
step $t_1$ and a key $(g_2,v_2)$ at step $t_2 > t_1$, then
$g_1 + v_1 \leq g_2 + v_2$, and if $g_1 + v_1 = g_2 + v_2$ then
$g_1 \leq g_2$.}
For the $f$-part, every state inserted into the open list during the
expansion of a bucket with key $(g,v)$ is a successor $s[o]$ of a state
$s$ with $h(s) = v$ and has key $(g + \cost(o),\, h(s[o]))$, whose
$f$-value satisfies
$g + \cost(o) + h(s[o]) \geq g + h(s) = g + v$ by consistency. Hence
expanding a minimum-$f$ key only inserts keys of at least that
$f$-value, the minimum $f$-value of the open list never decreases, and
selected $f$-values are non-decreasing over time. For the $g$-part,
suppose for contradiction that there are selections at steps
$t_1 < t_2$ of keys with equal $f$-value $F$ and $g_2 < g_1$, and among
all such violations choose one minimizing $t_2$. The bucket of
$(g_2,v_2)$ must be empty at step $t_1$: otherwise the tie-breaking
rule would select it instead of $(g_1,v_1)$. Its content is also not
inserted at step $t_1$ itself, since states inserted by the expansion
of $(g_1,v_1)$ have $g$-component $g_1 + \cost(o) > g_1 > g_2$. So
every state of the bucket is inserted at some step $t'$ with
$t_1 < t' < t_2$ by expanding a key $(g',v')$ with
$g' = g_2 - \cost(o) < g_2$ for some operator $o$. Consistency gives
$g' + v' \leq g_2 + v_2 = F$ as above, and the $f$-part applied to
$t_1 < t'$ gives $g' + v' \geq F$, so $g' + v' = F$ and $g' < g_1$:
the pair $(t_1, t')$ is a violation with $t' < t_2$, contradicting
minimality.

\emph{Claim 2 (no re-population): after a key $(g,v)$ is selected at
step $t$, no state with key $(g,v)$ is inserted into the open list;
consequently every key is selected at most once, and at its selection
its bucket contains every state that was ever generated with this key
and not expanded before.}
Suppose a state with key $(g,v)$ is inserted at a step $t' > t$ by
expanding a key $(g',v')$. Then $g' + v' \leq g + v$ by consistency as
in Claim~1, and $g' + v' \geq g + v$ by the $f$-part of Claim~1, so the
$f$-values are equal; but $g' = g - \cost(o) < g$ contradicts the
$g$-part. A second selection of the same key would require an
insertion after step $t$, which is excluded. Finally, a state generated
with key $(g,v)$ before step $t$ enters the bucket unless it is removed
by the closed-list test, which removes only previously expanded states,
and states are removed from the open list only when their bucket is
selected.

\emph{Claim 3 (optimal $g$-values): every state $s$ in a selected
bucket with key $(g,v)$ satisfies $\gstar(s) = g$ and $h(s) = v$.}
We use induction on the selection step; Claims~1 and~2 hold
unconditionally and may be used freely. That $h(s) = v$ holds by
construction of the level-set partition, and $\gstar(s) \leq g$ because
$s$ is reached by a path of cost $g$. Suppose $\gstar(s) < g$ for some
$s$ in the bucket selected at the current step, and fix a cheapest path
$p_0 = \sinit, \dots, p_k = s$, all of whose prefixes are cheapest
paths to their end states. Let $j$ be the smallest index such that
$p_j$ has not been expanded before the current step; $j$ exists because
$s$ qualifies, and $j \geq 1$ because \bddastar\ selects the initial
bucket $\{\sinit\}$ first. By the induction hypothesis, every earlier
expansion of $p_{j-1}$ has key $(\gstar(p_{j-1}), h(p_{j-1}))$, so
$p_j$ was generated with key $(\gstar(p_j), h(p_j))$ before the current
step. This key has not been selected yet: a selection after the
generation of $p_j$ would have expanded $p_j$ (Claim~2), contradicting
the choice of $j$, and a selection before the generation would
contradict Claim~2 directly. Hence the open list currently contains the
key $(\gstar(p_j), h(p_j))$, and consistency along the path segment
from $p_j$ to $s$ yields
$\gstar(p_j) + h(p_j) \leq \gstar(s) + h(s) < g + v$, using
$\gstar(s) < g$ and $h(s) = v$. This contradicts minimum-$f$ selection.

\emph{Claim 4 (completeness): every state $s$ with
$\gstar(s) = g < \copt$ and $g + h(s) \leq \copt$ is generated with key
$(g, h(s))$, and this key is selected before the search terminates.}
By strong induction on $g$. \emph{Generation:} for $g = 0$, $s = \sinit$
initializes the open list with key $(0, h(\sinit))$. For $g > 0$, let
$p$ be the predecessor of $s$ on a cheapest path, with
$\gstar(p) = g - \cost(o) < g$ and, by consistency,
$\gstar(p) + h(p) \leq \gstar(s) + h(s) \leq \copt$. By the induction
hypothesis, $p$ is generated with key $(\gstar(p), h(p))$ and this key
is selected before termination; by Claim~2 the bucket then contains
$p$, because any earlier expansion of $p$ would, by Claim~3, have
occurred with the same key, which is selected only once. Expanding the
bucket generates $s$ at cost $\gstar(p) + \cost(o) = g$ with key
$(g, h(s))$. \emph{Selection:} let $K = (g, h(s))$, with
$f(K) \leq \copt$ and $g < \copt$. From the step at which $K$ first
becomes nonempty, $K$ remains in the open list until it is selected
(Claim~2). The search terminates only when a selected bucket contains a
goal state, and by Claim~3 a goal state in a selected bucket has key
$(\gstar(s_G), 0)$ with $\gstar(s_G) \geq \copt$, since consistent
heuristics assign $0$ to goal states and \copt\ is the cheapest goal
cost. While $K$ is open, no such key is selected: its $f$-value is at
least $\copt \geq f(K)$, and in case of equality its $g$-component
$\gstar(s_G) \geq \copt > g$ loses the tie-break against $K$ or another
open key with smaller $g$-component. Hence the search does not
terminate while $K$ is open. Moreover, while $K$ is open every
selection picks a key with $f$-value at most $f(K) \leq \copt$; there
are at most $(\copt + 1)(\copt + 2)/2$ keys $(g', v')$ of nonnegative
integers with $g' + v' \leq \copt$, each selected at most once
(Claim~2), and the open list is nonempty while $K$ is open. Hence after
finitely many steps $K$ itself is selected, before termination.

\emph{Claim 5 (no other expansions): every bucket selected before the
terminating selection has a key $(g,v)$ with $g + v \leq \copt$ and
$g < \copt$.}
Fix a cheapest plan $q_0 = \sinit, \dots, q_r$ with $q_r$ a goal state
and all prefixes cheapest, so that consistency yields
$\gstar(q_l) + h(q_l) \leq \copt$ for every $l$. First we exhibit a
witness key: at every step at which the search has not yet terminated,
let $l$ be the smallest index such that $q_l$ is unexpanded, which
exists because $q_r$ is unexpanded before termination. If $l = 0$, the
initial key $(0, h(\sinit))$ has not been selected, since its bucket
contains $\sinit$, and is therefore open. If $l \geq 1$, then $q_{l-1}$
was expanded, by Claim~3 with key $(\gstar(q_{l-1}), h(q_{l-1}))$, so
$q_l$ was generated with key $K_l = (\gstar(q_l), h(q_l))$; exactly as
in the proof of Claim~3, $K_l$ has not been selected and is therefore
open. In all cases the open list contains a key with $f$-value at most
\copt, namely $(0, h(\sinit))$ or $K_l$; call it the witness. Second,
minimum-$f$ selection against the witness implies that every selected
key satisfies $g + v \leq \copt$. Third, a selected key with
$g + v \leq \copt$ and $g \geq \copt$ must equal $(\copt, 0)$, and we
show that its selection is the terminating one. Consider the step at
which $(\copt, 0)$ is selected and let $l$ be as above. If $l < r$,
the witness $K_l$ is open with $f$-value at most \copt\ and
$g$-component $\gstar(q_l) < \copt$; then either its $f$-value is
smaller than \copt, contradicting minimum-$f$ selection of
$(\copt, 0)$, or it equals \copt\ and the tie-breaking rule prefers
$K_l$ or another open key of smaller $g$-component over $(\copt, 0)$,
again a contradiction. Hence $l = r$, so $q_r$ was generated with key
$(\copt, 0)$ and, being unexpanded, lies in its bucket at selection
time by Claim~2. The selected bucket thus contains the goal state
$q_r$ and the search terminates. Consequently, every selection before
the terminating one has $g + v \leq \copt$ and $g < \copt$.

\emph{Assembly.} Let $(g,v)$ with $g + v \leq \copt$ and $g < \copt$
and $L_g \cap \level{h}{v} \neq \emptyset$. Every
$s \in L_g \cap \level{h}{v}$ satisfies the premise of Claim~4, so its
key $(g,v)$ is selected exactly once before termination (Claims~2
and~4), and at this selection the bucket contains $s$ (Claim~2, with
prior expansion of $s$ excluded because, by Claim~3, it would have
occurred at the same, uniquely selected key). Conversely, by Claim~3
the selected bucket contains only states with $\gstar = g$ and
$h$-value $v$. Hence the expanded BDD represents exactly
$L_g \cap \level{h}{v}$. Keys whose intersection is empty are never
selected, since buckets enter the open list only when populated, and by
Claim~5 no key outside the stated range is expanded before the search
stops.
\end{proof}

With the bucket identity in place, the effort bound follows from the
partition bound applied layer by layer.

\begin{theorem}[effort bound]
\label{thm-effort}
Let $h$ be consistent with $\width_\pi(h) \leq W$. Then
\[
\effort(\task, h)
\;\leq\; W^{2} n \cdot \bigl(\effort(\task, \hblind) + 2\copt\bigr).
\]
\end{theorem}

\begin{proof}
By Lemma~\ref{lem-identity} and Definition~\ref{def-effort},
\[
\effort(\task,h)
= \sum_{g < \copt} \;\sum_{\substack{v \in \img(h)\\ g + v \leq \copt}}
\bddsize{L_g \cap \level{h}{v}}
\;\leq\; \sum_{\substack{g < \copt\\ L_g \neq \emptyset}}
\;\sum_{v \in \img(h)}
\bddsize{L_g \cap \level{h}{v}},
\]
where the restriction to nonempty layers loses nothing because all
buckets of an empty layer are empty and the BDD of the empty set has no
inner nodes. By Theorem~\ref{thm-partition}, the inner sum is at most
$W^{2} n \,(\bddsize{L_g} + 2)$ for every $g$. Writing $N$ for the
number of nonempty layers $L_g$ with $g < \copt$ and using
$\effort(\task,\hblind) = \sum_{g < \copt} \bddsize{L_g}$, we obtain
\begin{equation}
\effort(\task,h)
\;\leq\; W^{2} n \,\bigl(\effort(\task,\hblind) + 2N\bigr),
\label{eq-effort-nonempty}
\end{equation}
and $N \leq \copt$ yields the claim.
\end{proof}

\begin{corollary}
\label{cor-effort}
If additionally every nonempty layer $L_g$ with $g < \copt$ is a proper
subset of the set of all Boolean assignments, then
$\effort(\task, h) \leq 3\, W^{2} n \cdot \effort(\task, \hblind)$.
\end{corollary}

\begin{proof}
The BDD of a nonempty proper subset has at least one inner node, so
$N \leq \effort(\task,\hblind)$ in
Equation~\ref{eq-effort-nonempty}.
\end{proof}

The condition of Corollary~\ref{cor-effort} fails only in degenerate
tasks. Theorem~\ref{thm-effort} extends beyond \bddastar:
\citet{speck-et-al-icaps2020} observe that Lazy \bddastar\
\citep{torralba-phd2015} and SetA\ensuremath{^{*}}
\citep{jensen-et-al-aij2008}, and by the same argument its
cost-general variant \ghseta, expand exactly the same state sets as
\bddastar, differing only in when and how the partition is computed, so
the bound on the cumulative size of expanded BDDs applies to these
variants as well. Two honest limitations. First, the guarantee is
one-sided by design: it caps the overhead of using $h$, while the
potential savings, namely the layers and dead ends that $h$ prunes via
the constraints $g + v \leq \copt$ and $v \neq \infty$, remain
unbounded, which is the desired asymmetry. Second, the baseline is blind
\emph{forward} search; the strongest blind configuration in practice is
bidirectional \citep{torralba-et-al-aij2017,speck-et-al-jair2025}, and
extending the analysis to bidirectional \bddastar, where the backward
frontier interacts with forward buckets, is future work.

\subsection{Exponential Width Is Necessary for the Known Pathology}

We now revisit the family $\task_n$ of \citet{speck-et-al-icaps2020} on
which \bddastar\ with \hstar\ is exponentially worse than with \hblind.
The task has binary variables
$v_1, \dots, v_{2n}$ and $x_0, \dots, x_{2n+1}$; initially only $x_0$
is true, operators $o_i$ and $\bar o_i$ with precondition
$x_{i-1} \wedge \neg x_i$ set $x_i$ true while setting $v_i$ true or
false, respectively, and once $x_0, \dots, x_{2n}$ are all true, a goal
operator $g_i$ with precondition $x_{2n} \wedge v_i \wedge v_{n+i}$ for
some $i \in \{1, \dots, n\}$ achieves the goal
$x_{2n+1}$. The variable order is
$\pi_n = x_0, \dots, x_{2n+1}, v_1, \dots, v_{2n}$, and the exponential
blowup arises in the bucket of non-dead-end states at $g = 2n$, whose
characteristic function contains
$\bigvee_{i=1}^{n} (v_i \wedge v_{n+i})$, a function with exponential
OBDD size under $\pi_n$ \citep{kissmann-phd2012,speck-et-al-icaps2020}.

\begin{theorem}[width lower bound]
\label{thm-lb}
For the tasks $\task_n$ and orders $\pi_n$ of
\citet{speck-et-al-icaps2020},
$\width_{\pi_n}(\hstar) \geq 2^{n}$.
\end{theorem}

\begin{proof}
We exhibit $2^n$ prefixes with pairwise distinct \hstar-cofactors. For
$A \subseteq \{1, \dots, n\}$ let $\rho_A$ be the prefix that assigns
the $x$-variables the encoding of chain position $2n$, i.e.,
$x_0, \dots, x_{2n}$ true and $x_{2n+1}$ false, and assigns $v_i = 1$
iff $i \in A$ for $i \in \{1, \dots, n\}$. All $\rho_A$ fix the same
variables $x_0, \dots, x_{2n+1}, v_1, \dots, v_n$, so their cofactors
lie in $\sub{m}{\hstar}$ for the same level $m = (2n + 2) + n$. Let
$A \neq A'$ and pick without loss of generality
$i \in A \setminus A'$. Define the suffix $\sigma$ over
$v_{n+1}, \dots, v_{2n}$ by $v_{n+i} = 1$ and $v_{n+j} = 0$ for
$j \neq i$. The state $\rho_A \sigma$ satisfies
$v_i \wedge v_{n+i}$, so a goal operator is applicable and
$\hstar(\rho_A \sigma) = 1$. The state $\rho_{A'} \sigma$ satisfies
$v_j \wedge v_{n+j}$ for no $j$: for $j = i$ because $v_i = 0$ under
$\rho_{A'}$, and for $j \neq i$ because $v_{n+j} = 0$ under $\sigma$.
No other operator is applicable in either state, since every $o_i$ and
$\bar o_i$ requires $\neg x_i$ and $x_1, \dots, x_{2n}$ are all true.
Hence $\rho_{A'} \sigma$ is a dead
end and $\hstar(\rho_{A'} \sigma) = \infty \neq 1 =
\hstar(\rho_A \sigma)$, so
$\hstar|_{\rho_A} \neq \hstar|_{\rho_{A'}}$, and
$|\sub{m}{\hstar}| \geq 2^{n}$. Since the argument evaluates \hstar\
only on states, it applies to every total extension of \hstar\ to
$\{0,1\}^n$ (Remark~\ref{rem-encoding}).
\end{proof}

Theorems~\ref{thm-effort} and~\ref{thm-lb} together delimit the
phenomenon discovered by \citet{speck-et-al-icaps2020}: any heuristic
that deteriorates \bddastar\ by a superpolynomial factor over blind
search must have superpolynomial width under the search order, and the
perfect heuristic on their family indeed has exponential width.
Width-boundedness is a sufficient condition for safety, not a necessary
one: a heuristic of large width may still happen to induce small
buckets. Note also that \hstar\ on $\task_n$ takes at most
$2n + 3$ distinct values, so a small value range does not imply small
width; the width captures how much prefix information the level sets
entangle, not how many there are.

% ============================ instances ==================================
\section{Width-Bounded Heuristic Families}
\label{sec-instances}

We now show that several practically important heuristic families are
width-bounded, with widths governed by natural parameters. Throughout,
$\pi$ is variable-contiguous, $d = \max_{v \in \vars} |D_v|$ and
$m = |\vars|$, so $n \leq m \lceil \log_2 d \rceil$.

\subsection{Potential Heuristics}

A potential heuristic \citep{pommerening-et-al-aaai2015} assigns a
potential $P(v, w) \in \mathbb{R}$ to every fact and evaluates to
$\hpot(s) = \sum_{v \in \vars} P(v, s(v))$. Admissibility and
consistency are enforced by linear constraints.
\citet{fiser-et-al-aaai2022} round potentials to integers inside the
optimization to integrate them into \ghseta; we adopt their integer
setting.

\begin{proposition}
\label{prop-pot}
Let \hpot\ have integer potentials with $|P(v,w)| \leq M$ for all
facts, extended to non-state assignments by the same sum over the bit
encoding's induced table, with table values on invalid bit codes also
in $[-M, M]$, e.g., $0$, and let $\pi$ be variable-contiguous. Then
$\width_\pi(\hpot) \leq d \, (2mM + 1)$ and
$|\img(\hpot)| \leq 2mM + 1$.
\end{proposition}

\begin{proof}
Consider a prefix $\rho$ that fixes the variables
$v_1, \dots, v_j$ completely and the first $b < \lceil \log_2 d \rceil$
bits of $v_{j+1}$. The cofactor $\hpot|_\rho$ is determined by the pair
consisting of the accumulated sum
$\sum_{j' \leq j} P(v_{j'}, \rho(v_{j'}))$, an integer in
$[-mM, mM]$, and the assignment to the $b$ read bits of $v_{j+1}$,
because the remaining function is the fixed table sum over the
unread variables shifted by the accumulated constant and restricted to
the read bits. There are at most $2mM + 1$ accumulated sums and at most
$2^{b} \leq d$ bit prefixes, so
$|\sub{i}{\hpot}| \leq d\,(2mM + 1)$ for every level $i$. The value
range consists of total sums, integers in $[-mM, mM]$, at most
$2mM + 1$ many; shifting by a constant to make values nonnegative
changes neither count.
\end{proof}

\begin{corollary}
\label{cor-pot}
For consistent integer potential heuristics with magnitudes at most
$M$, under the conditions of Corollary~\ref{cor-effort},
$\effort(\task, \hpot) \leq
3\, d^{2} (2mM + 1)^{2}\, n \cdot \effort(\task, \hblind)$,
a factor polynomial in the task size and $M$.
\end{corollary}

Corollary~\ref{cor-pot} is, to our knowledge, the first formal
performance guarantee for the operator-potential approach of
\citet{fiser-et-al-aaai2022,fiser-et-al-aij2024}, whose reported
behavior---informative guidance with rarely increased BDD
sizes---matches exactly what a small width predicts. It also explains
why their integer rounding matters twice: rounding coarsens the value
range \emph{and} caps the width through $M$.

\subsection{Pattern Databases over Pattern Prefixes}

A pattern database \citep[PDB;][]{culberson-schaeffer-compint1998,%
edelkamp-ecp2001} projects \task\ onto a pattern
$P \subseteq \vars$ and stores the abstract goal distance
$h^{P}(s)$ of the abstract state induced by $s$.

\begin{proposition}
\label{prop-pdb}
Let $P \subseteq \vars$ such that the variables of $P$ form a prefix of
the variable-contiguous order $\pi$, and let
$|\states^{P}| = \prod_{v \in P} |D_v|$ denote the number of abstract
states. Then $\width_\pi(h^{P}) \leq d \cdot |\states^{P}|$.
\end{proposition}

\begin{proof}
For a prefix $\rho$ inside the pattern block, the cofactor
$h^{P}|_\rho$ is determined by the assignment to the read pattern
variables together with the read bits of the current variable: two
prefixes inducing the same partial abstract state and bit buffer have
identical completions. The number of partial abstract states is at most
$|\states^{P}|$ and the bit buffer contributes a factor at most $d$ as
in Proposition~\ref{prop-pot}. For prefixes beyond the pattern block,
$h^{P}|_\rho$ is the constant $h^{P}$-value of the abstract state fixed
by $\rho$, so the cofactor count is at most
$|\img(h^{P})| \leq |\states^{P}|$.
\end{proof}

The bound $d \cdot |\states^{P}|$ is coarse; the exact width is the
maximal level size of the PDB's quasi-reduced ADD, which reduction
often makes much smaller. The requirement that $P$ be a prefix of $\pi$
is the price of alignment between abstraction and representation:
patterns scattered through the order can entangle prefix information
much like Theorem~\ref{thm-lb}. By Proposition~\ref{prop-closure},
maximizing or admissibly adding a few prefix PDBs, e.g., under a cost
partitioning, multiplies the width bounds.

\subsection{Linear Merge-and-Shrink}

Merge-and-shrink (M\&S) abstractions
\citep{helmert-et-al-jacm2014} iteratively merge factors of a factored
transition system and shrink intermediate abstractions to at most $N$
abstract states. With a \emph{linear} merge strategy, the abstraction
mapping is representable by cascading tables, equivalently a factored
mapping \citep{helmert-et-al-icaps2015}: a chain of tables
$\tau_1, \dots, \tau_m$ where $\tau_j$ maps the pair of the previous
intermediate abstract state and the value of variable $v_j$ to the next
intermediate abstract state.

\begin{proposition}
\label{prop-ms}
Let $h^{\alpha}$ be an M\&S heuristic built with a linear merge
strategy whose variable order agrees with the variable-contiguous
$\pi$ and with intermediate abstraction sizes at most $N$. Then
$\width_\pi(h^{\alpha}) \leq d \cdot N$.
\end{proposition}

\begin{proof}
For a prefix $\rho$ fixing $v_1, \dots, v_j$ and $b$ bits of
$v_{j+1}$, the cofactor $h^{\alpha}|_\rho$ is determined by the
intermediate abstract state $\tau_j(\cdots)$ reached after processing
$v_1, \dots, v_j$ through the cascading tables, of which there are at
most $N$, together with the $b$-bit buffer, at most $d$ cases: the
remaining computation applies the fixed tables
$\tau_{j+1}, \dots, \tau_m$ and the final distance table to the suffix.
\end{proof}

Proposition~\ref{prop-ms} recasts a known compactness result in a new
role. Running \bddastar\ with M\&S heuristics compiled into BDDs goes
back to \citet{edelkamp-et-al-ecai2012}, and \citet{torralba-phd2015}
proved that linear M\&S heuristics have
BDD representations with at most polynomial overhead, enabling their
use in symbolic search; Proposition~\ref{prop-ms} sharpens the
representational statement to a per-level width bound and, through
Theorems~\ref{thm-partition} and~\ref{thm-effort}, converts it for the
first time into a guarantee on \emph{search behavior}: the shrink size
limit $N$, imposed anyway to keep M\&S construction tractable, doubles
as a bound on the harm the heuristic can do to the symbolic search that
uses it. Our merge-and-shrink configuration is thus deliberately the
configuration of \citet{edelkamp-et-al-ecai2012} made safe: same
algorithmic skeleton, now with an alignment condition (merge order
matching the search's variable order), a width guarantee, and---in
Section~\ref{sec-experiments}---an account of when it pays
off.\footnote{Our default configuration merges along a causal-graph
order that need not coincide with the search's Gamer order; the
statistics we report therefore measure the width of the constructed ADD
directly \emph{in the search order}, which is the quantity the theorems
consume. Section~\ref{sec-experiments} evaluates enforcing the
alignment condition.} In contrast, no small-width representation is known for
$h^{\max}$, $h^{\mathrm{add}}$ or LM-cut, whose values depend on
optimization problems that entangle all variables, and
Theorem~\ref{thm-lb} shows that perfect goal-distance information can
force exponential width.

% ============================ synthesis ==================================
\section{Synthesizing Width-Bounded Heuristics}
\label{sec-synthesis}

The results above invert the usual heuristic-selection problem for
symbolic search: instead of choosing the most informative heuristic and
hoping that the partition stays benign, we propose to optimize
informativeness \emph{subject to a width budget}. For potential
heuristics this is directly actionable. The LP of
\citet{pommerening-et-al-aaai2015} maximizes, e.g., the heuristic value
of the initial state or the average value over samples, subject to
admissibility and consistency constraints. Following
\citet{fiser-et-al-aaai2022} we require integer potentials, and we add
box constraints $|P(v,w)| \leq M$, which by
Proposition~\ref{prop-pot} enforce
$\width_\pi(\hpot) \leq d\,(2mM+1)$. Coarsening potentials to multiples
of an integer $k$ tightens the width bound to $d\,(2mM/k + 1)$ at a
bounded loss of accuracy, since down-rounding sums to multiples of $k$
preserves admissibility. The budget $M$ (or $M/k$) is thus a knob that
interpolates between blind search ($M = 0$) and unconstrained potential
heuristics, with Corollary~\ref{cor-pot} guaranteeing the overhead at
every setting. A sharper objective minimizes the number of distinct
reachable partial sums directly, which is the exact width up to the bit
buffer; we formulate this as a MIP but expect the box-constraint
relaxation to be the practical variant. For prefix PDBs, synthesis
amounts to pattern selection constrained to prefixes of $\pi$, which
composes with the causal-graph-based orders of
\citet{kissmann-edelkamp-aaai2011} and dovetails with the variable
relaxation strategy used for symbolic perimeter abstractions
\citep{torralba-et-al-aij2018}.

% ============================ experiments ================================
\section{Experiments}
\label{sec-experiments}

We implemented width instrumentation and all three width-bounded
heuristic families---width-capped integer potential heuristics, prefix
pattern databases and linear merge-and-shrink heuristics---in the SymK
planning system \citep{speck-et-al-jair2025}, which extends Fast Downward
\citep{helmert-jair2006}, and used the Downward Lab toolkit
\citep{seipp-et-al-zenodo2017} for running experiments on Intel Xeon Gold
6130 processors at 2.1\,GHz. Our benchmark set consists of all tasks
without conditional effects and axioms from the optimal sequential tracks
of the IPCs 1998--2023, a subset of the suite of
\citet{fiser-et-al-aij2024}. Because the width-bounded search assumes
positive operator costs (Section~\ref{sec-effort}), the heuristic
configurations are undefined on the 320 tasks with zero-cost operators;
we therefore report coverage on the remaining 1377 tasks, which all
configurations attempt. Two exceptions to this base are flagged where
they occur: the bidirectional pruning study uses the full suite of 1697
tasks, since pruning-only search needs no positive-cost restriction, and
repeated runs of bidirectional SymK differ by one or two solved tasks
because its direction selection is timing-based; we always compare runs
from the same sweep. We limit time to 30\,minutes and memory to
8\,GiB. All code, benchmarks and experiment data are available in a
supplementary archive.\footnote{A Zenodo DOI will be added for the
camera-ready version.}

We focus on coverage (number of solved tasks), the expansion size
(Definition~\ref{def-effort}) measured in BDD nodes, and the
\emph{fragmentation ratio}, i.e., the total size of the buckets of a
layer divided by the size of the unpartitioned layer BDD, which
Theorem~\ref{thm-partition} predicts to be at most polynomial in the
measured width. We compare against forward and bidirectional blind
symbolic search in SymK, the operator-potential planner of
\citet{fiser-et-al-aij2024}, \symba\
\citep{torralba-et-al-aij2017} and the explicit-state cost-partitioning
planner Scorpion \citep{seipp-et-al-jair2020}. Our evaluation addresses
four questions. (Q1)~Does the measured fragmentation ratio respect the
width bound? (Q2)~How do coverage and expansion size change with the
width budget $M$? (Q3)~How do the width-bounded families compare to blind
search and to the state of the art? (Q4)~Does the width bound behave as a
usable accuracy--overhead knob?

\paragraph{Q1: Fragmentation.}
Across all solved tasks and all three families, partitioning a layer by
heuristic value increases its total BDD size only marginally: the median
fragmentation ratio is $1.17$ for prefix PDBs, $1.29$ for potentials and
$1.31$ for linear merge-and-shrink, and it never exceeds $2.7$, even
though the measured width upper bound $U = A + \values$
(Proposition~\ref{prop-add}) ranges over five orders of magnitude across
domains (medians $141$, $576$ and $1171$, maxima up to $5{\cdot}10^{6}$).
We emphasize the logical status of these statistics: $U$ is computed
from the heuristic's ADD \emph{in the search's variable order}, which is
exactly the quantity Theorems~\ref{thm-bucket}--\ref{thm-effort}
consume, so the guarantees cover every reported run; the closed-form
parameter bounds of Section~\ref{sec-instances} (such as $dN$ for
merge-and-shrink) additionally require the stated alignment conditions,
which our implementation does not enforce for the merge order.
The fragmentation ratio thus stays small and far below the worst case of
Theorem~\ref{thm-partition}, confirming that width-bounded partitioning
does not blow up BDD sizes in practice. As a sanity anchor we also ran
blind forward search on the family $\task_n$ of
\citet{speck-et-al-icaps2020}: its layers stay linear in $n$ (the largest
expanded layer has $18$, $34$ and $50$ BDD nodes for $n = 4, 8, 12$),
whereas the perfect heuristic has width $2^{n}$ there
(Theorem~\ref{thm-lb}), the very blow-up that width-boundedness rules out.

\begin{table}[t]
\centering
\begin{tabular}{lr}
\toprule
Configuration & Coverage \\
\midrule
Blind forward (SymK) & 682 \\
Potentials ($M = 8$), \bddastar & 622 \\
Prefix PDB, \bddastar & 685 \\
Linear merge-and-shrink, \bddastar & 704 \\
Potentials ($M = 8$), pruning-only & 683 \\
Prefix PDB, pruning-only & 682 \\
Linear merge-and-shrink, pruning-only & 681 \\
\midrule
Blind bidirectional (SymK) & 777 \\
A+I operator potentials \citep{fiser-et-al-aij2024} & 652 \\
\symba\ \citep{torralba-et-al-aij2017} & 724 \\
Scorpion \citep{seipp-et-al-jair2020} & \textbf{912} \\
\bottomrule
\end{tabular}
\caption{Coverage on the 1377 positive-cost tasks. The upper block lists
blind forward search and our three width-bounded families, each in the
\bddastar\ (partitioning) and the pruning-only variant; the lower block
lists external baselines.}
\label{tab-coverage}
\end{table}

\paragraph{Q3: Coverage.}
Table~\ref{tab-coverage} reports total coverage and
Table~\ref{tab-domains} the per-domain breakdown. All three width-bounded
families are competitive with blind forward search: prefix PDBs ($685$)
and linear merge-and-shrink ($704$) solve slightly more tasks, capped
potentials ($622$) slightly fewer. Since Theorem~\ref{thm-effort} only
bounds the \emph{overhead} of adding a heuristic, matching blind coverage
while enjoying the width guarantee is the expected outcome; the
merge-and-shrink family, whose abstractions capture more global structure
than a single potential function, improves on it, and both improve on the
A+I operator-potential planner ($652$). A caveat on the latter
comparison: we run the A+I planner in its default single-potential
configuration (forward and backward operator potentials from one LP with
the initial-state objective); the strongest configurations reported by
\citet{fiser-et-al-aij2024} combine several potential functions, so
$652$ should be read as a lower bound on what that system can achieve on
this suite. The remaining configurations that
solve more tasks all use search paradigms outside the scope of our
forward-search guarantee (Theorem~\ref{thm-effort}): blind
\emph{bidirectional} SymK ($777$) and \symba\ ($724$) are bidirectional,
and Scorpion ($912$) is explicit-state. Closing the gap to bidirectional
search---where the backward frontier interacts with forward buckets---is
future work. The result we do claim
is the one the theory predicts: width-bounded heuristics can be added to
forward symbolic search essentially for free, and yield modest coverage
gains.


Why do safe potentials nevertheless lose coverage? The guarantee bounds
the \emph{size} of the expanded BDDs, but expanding up to $\values$
buckets per layer multiplies the number of image operations, each of
which carries overhead beyond its argument's size; potentials have many
values (median $53$ on this suite) and merge-and-shrink few (median
$16$), which matches the observed coverage split. A cost model that
makes this operational price explicit, and search variants that avoid
it, are the subject of ongoing work.
\paragraph{Merge-order alignment.}
Proposition~\ref{prop-ms} requires the merge order to match the search's
variable order for the closed-form $dN$ bound; our default
merge-and-shrink configuration instead uses a causal-graph order and is
covered by the theorems through its measured width. Enforcing the
alignment condition does what the theory predicts for the
representation---the measured width drops (median $1380$ to $1091$, and
the $90$th percentile falls fivefold, from $50\,059$ to
$9\,751$)---but costs coverage ($687$ vs.\ $704$; five tasks won,
twenty-two lost): the causal-graph merge order produces more informative
abstractions under the same shrink limit, and here accuracy outweighs
representational compactness. Width, in other words, is a currency, not
the objective---consistent with our one-sided guarantee, which caps the
harm of a heuristic but says nothing about its benefit.

\paragraph{Q2 and Q4: The width knob.}
Table~\ref{tab-knob} varies the potential magnitude cap $M$. At $M = 0$
the heuristic is uniformly zero and the search reduces to blind forward
search: the expansion size matches it to within $0.6\%$, confirming the
blind-equivalence of our implementation. As $M$ grows, the geometric-mean
expansion-size overhead over blind forward search stays a small bounded
factor---at most $1.9$ and shrinking to $1.3$---never approaching the
exponential deterioration possible for heuristics of unbounded width
(Theorem~\ref{thm-lb}). Coverage is highest around $M = 8$, which we fix
as the default. The cap thus behaves as intended: a knob that trades
heuristic accuracy against a guaranteed, and empirically small, limit on
representational overhead.

\begin{table}[t]
\centering
\begin{tabular}{lrrrrrrr}
\toprule
$M$ & 0 & 1 & 2 & 4 & 8 & 16 & $\infty$ \\
\midrule
Coverage & 680 & 601 & 602 & 608 & 622 & 623 & 596 \\
Effort ratio & 1.01 & 1.92 & 1.62 & 1.45 & 1.33 & 1.33 & 1.28 \\
\bottomrule
\end{tabular}
\caption{Coverage and geometric-mean expansion-size ratio over blind
forward search for capped potentials with magnitude cap $M$
($\infty$ denotes an effectively unbounded cap).}
\label{tab-knob}
\end{table}

\begin{table}[p]
\centering
\small
\setlength{\tabcolsep}{4pt}
\begin{tabular}{lrrrrrrrrr}
\toprule
Domain & \# & b-fw & pot & PDB & M\&S & b-bd & A+I & SymBA & Scrp \\
\midrule
airport & 50 & 23 & 22 & 23 & 23 & 24 & 27 & 27 & 37 \\
barman-opt11 & 20 & 8 & 4 & 8 & 8 & 11 & 8 & 12 & 8 \\
barman-opt14 & 14 & 3 & 3 & 3 & 3 & 6 & 3 & 6 & 3 \\
blocks & 35 & 21 & 21 & 21 & 28 & 31 & 27 & 33 & 28 \\
childsnack-opt14 & 20 & 4 & 2 & 4 & 2 & 3 & 2 & 2 & 0 \\
depot & 22 & 5 & 4 & 6 & 8 & 7 & 7 & 7 & 13 \\
driverlog & 20 & 11 & 11 & 11 & 13 & 14 & 12 & 14 & 15 \\
floortile-opt11 & 20 & 14 & 14 & 14 & 14 & 14 & 14 & 14 & 6 \\
floortile-opt14 & 20 & 20 & 17 & 20 & 17 & 20 & 20 & 20 & 5 \\
freecell & 80 & 20 & 22 & 20 & 20 & 25 & 47 & 26 & 68 \\
grid & 5 & 2 & 1 & 2 & 2 & 2 & 2 & 2 & 3 \\
gripper & 20 & 20 & 20 & 20 & 20 & 20 & 20 & 20 & 7 \\
hiking-opt14 & 20 & 16 & 15 & 16 & 16 & 18 & 13 & 18 & 19 \\
logistics00 & 28 & 16 & 20 & 17 & 21 & 20 & 0 & 0 & 28 \\
logistics98 & 35 & 5 & 5 & 5 & 6 & 6 & 5 & 5 & 12 \\
miconic & 150 & 102 & 88 & 100 & 79 & 117 & 78 & 114 & 145 \\
movie & 30 & 30 & 30 & 30 & 30 & 30 & 30 & 30 & 30 \\
mprime & 35 & 25 & 11 & 25 & 27 & 26 & 27 & 27 & 33 \\
mystery & 30 & 15 & 8 & 15 & 17 & 15 & 15 & 14 & 19 \\
nomystery-opt11 & 20 & 13 & 13 & 13 & 16 & 16 & 14 & 16 & 20 \\
organic-synthesis-opt18 & 20 & 7 & 7 & 7 & 7 & 7 & 13 & 13 & 7 \\
parking-opt11 & 20 & 0 & 0 & 0 & 4 & 1 & 1 & 2 & 8 \\
parking-opt14 & 20 & 0 & 0 & 0 & 4 & 1 & 3 & 4 & 8 \\
pipesworld-notankage & 50 & 16 & 18 & 17 & 21 & 17 & 19 & 16 & 27 \\
pipesworld-tankage & 50 & 17 & 18 & 17 & 18 & 17 & 17 & 15 & 18 \\
psr-small & 50 & 50 & 50 & 50 & 50 & 50 & 50 & 50 & 50 \\
rovers & 40 & 14 & 11 & 14 & 14 & 14 & 12 & 14 & 13 \\
satellite & 36 & 7 & 7 & 7 & 8 & 12 & 7 & 12 & 12 \\
scanalyzer-08 & 30 & 12 & 13 & 12 & 13 & 12 & 12 & 12 & 19 \\
scanalyzer-opt11 & 20 & 9 & 10 & 9 & 10 & 9 & 10 & 9 & 16 \\
snake-opt18 & 20 & 7 & 4 & 7 & 9 & 7 & 0 & 0 & 13 \\
storage & 30 & 15 & 14 & 15 & 15 & 15 & 0 & 0 & 16 \\
termes-opt18 & 20 & 12 & 12 & 14 & 12 & 18 & 13 & 18 & 7 \\
tetris-opt14 & 17 & 10 & 10 & 10 & 10 & 11 & 10 & 11 & 11 \\
tidybot-opt11 & 20 & 15 & 16 & 16 & 17 & 16 & 10 & 13 & 17 \\
tidybot-opt14 & 20 & 9 & 10 & 11 & 12 & 12 & 2 & 5 & 13 \\
tpp & 30 & 8 & 8 & 8 & 8 & 8 & 12 & 9 & 11 \\
transport-opt08 & 30 & 11 & 11 & 11 & 11 & 14 & 11 & 14 & 15 \\
transport-opt11 & 20 & 8 & 6 & 7 & 6 & 10 & 6 & 11 & 12 \\
transport-opt14 & 20 & 7 & 6 & 7 & 7 & 9 & 5 & 9 & 10 \\
trucks & 30 & 11 & 11 & 11 & 11 & 13 & 14 & 14 & 17 \\
visitall-opt11 & 20 & 12 & 13 & 12 & 13 & 12 & 12 & 12 & 17 \\
visitall-opt14 & 20 & 6 & 9 & 6 & 7 & 7 & 6 & 6 & 13 \\
woodworking-opt08 & 30 & 21 & 12 & 20 & 21 & 28 & 21 & 28 & 30 \\
woodworking-opt11 & 20 & 16 & 7 & 14 & 14 & 20 & 15 & 20 & 20 \\
zenotravel & 20 & 9 & 8 & 10 & 12 & 12 & 0 & 0 & 13 \\
\midrule
Total & 1377 & 682 & 622 & 685 & 704 & 777 & 652 & 724 & 912 \\
\bottomrule
\end{tabular}
\caption{Per-domain coverage on the positive-cost tasks: blind forward
(b-fw), the three \bddastar\ families with bounded cofactor width
(potentials $M=8$, prefix PDB, linear merge-and-shrink), blind
bidirectional (b-bd), and the external baselines A+I, \symba\ and
Scorpion (Scrp).}
\label{tab-domains}
\end{table}

% =========================== related work ================================
\section{Related Work}

Our work connects three lines of research. First, analyses of heuristic
symbolic search: \citet{jensen-et-al-aij2008} identified the arithmetic
partitioning of frontiers as a bottleneck and mitigated it through
transition-relation branching partitions;
\citet{speck-et-al-icaps2020} proved that no heuristic, not even
\hstar, is guaranteed to help \bddastar; we give the matching sufficient
condition and show their pathology requires exponential width. Second,
compact representations of heuristics.
Our merge-and-shrink configuration deserves an explicit comparison,
since its algorithmic skeleton is not new:
\citet{edelkamp-et-al-ecai2012} already compiled PDB and M\&S heuristics
into vectors of BDDs and ran \bddastar, and \citet{torralba-phd2015}
proved polynomial-size BDD representations for linear M\&S and analyzed
ADD-to-BDD conversions; the factored mappings of
\citet{helmert-et-al-icaps2015}, whose cascading tables our level-set
construction walks, made the representation explicit. What none of these
works provides---and what, after the negative result of
\citet{speck-et-al-icaps2020}, made the whole approach appear unsafe---is
a bound on the effect of the representation on the \emph{search
frontier}. Our contribution to this line is therefore not the encoding
but the account of when it works: the alignment condition between merge
order and search variable order, the width bound
$\width_\pi(h^\alpha) \leq dN$ that turns the standard shrink limit $N$
into a safety knob (Proposition~\ref{prop-ms}), and the empirical
account (Section~\ref{sec-experiments}) of when the configuration pays
off---this family has few heuristic values, which keeps the operational
cost of bucket-wise expansion low. The
operator-potential line
\citep{fiser-et-al-aaai2022,fiser-et-al-aij2024} demonstrated
empirically that potential heuristics integrate well into \ghseta;
Corollary~\ref{cor-pot} supplies the missing guarantee and identifies
the responsible parameter. Third, performance guarantees for symbolic
search: \citet{speck-helmert-ecai2025} bound the overhead of symbolic
over explicit blind search via conjunctive transition-relation
partitioning; our results are complementary, bounding the overhead of
heuristic over blind \emph{symbolic} search via a property of the
heuristic, and combining the two guarantee regimes is an appealing
direction. Related but orthogonal are symbolic perimeter abstraction
heuristics \citep{torralba-et-al-ijcai2016a,torralba-et-al-aij2018},
which defer heuristic computation until blind search stalls, variable
ordering analyses
\citep{kissmann-hoffmann-jair2014} and EVMDD-based representations of
state-dependent cost structure \citep{speck-et-al-icaps2018}.

% ============================ conclusions ================================
\section{Conclusions}

We introduced heuristics of bounded cofactor width and proved that
bounded width is a sufficient condition for safety in symbolic search:
partitioning any state set by heuristic values costs at most a
polynomial factor in the width, \bddastar\ with a consistent width-$W$
heuristic incurs at most an $O(W^2 n)$ overhead over blind symbolic
forward search, and the known exponential deterioration requires
exponential width. Potential heuristics, prefix PDBs and linear
merge-and-shrink instantiate the class with widths controlled by the
potential magnitude, the abstract state count and the shrink limit,
giving in particular the first formal guarantee for the
operator-potential approach. The experiments matched the theory:
fragmentation is negligible wherever width is bounded, the magnitude cap
behaves as a knob with bounded overhead, and the merge-and-shrink family
converts its guarantee into a coverage improvement over blind forward
search. The alignment study adds the cautionary half of the message:
enforcing the condition that makes the closed-form width bound
applicable shrinks the representation but costs accuracy---width is a
currency to be spent, not an objective to be minimized. We see two open
directions: extending the guarantee to bidirectional symbolic search,
and understanding the operational cost of bucket-wise expansion, which
our coverage results show is not captured by node counts alone.

% ------------------------------------------------------------------------
% Bibliography from the shared aibasel/bib repo (make update-bib).
\bibliographystyle{plainnat}
\bibliography{bib/abbrv,bib/literatur,bib/crossref}
\end{document}